\documentclass[11pt]{article}
\usepackage[margin=1in]{geometry}
\usepackage{amsmath,amssymb,amsthm,mathtools,bm}
\usepackage{microtype}
\usepackage[T1]{fontenc}
\usepackage{lmodern}
\usepackage{booktabs,array,enumitem,caption}
\usepackage[hidelinks]{hyperref}
\setlist{nosep}
\newtheorem{theorem}{Theorem}[section]
\newtheorem{corollary}[theorem]{Corollary}
\newtheorem{lemma}[theorem]{Lemma}
\newtheorem{proposition}[theorem]{Proposition}
\theoremstyle{definition}
\newtheorem{definition}[theorem]{Definition}
\newtheorem{assumption}[theorem]{Assumption}
\newtheorem{model}[theorem]{Modeling Assumption}
\theoremstyle{remark}
\newtheorem{remark}[theorem]{Remark}

\newcommand{\dd}{\mathrm{d}}
\newcommand{\R}{\mathbb{R}}
\newcommand{\GM}{GM}
\newcommand{\FH}{F_H}
\newcommand{\Tr}{\operatorname{Tr}}

\title{\textbf{Nonsingular Gravitational Collapse via Hysterical Instability}\\[0.35em]
\large Hysterical Response Theory proves a stable finite-radius remnant}
\author{Andrew Bruce Boyd}
\date{Working mathematical formulation --- September 2026}

\begin{document}
\maketitle
\noindent\textit{Computational note.}
Proofs, exploratory experiments, and paper writing were assisted by generative AI.
Mathematical theorems stated in this paper are formally verified in Lean~4 in the
companion specifications, as faithful ports of the TeX arguments at the abstractions
recorded there, except results explicitly tagged as \emph{literature hypotheses}
(standard theorems cited from the literature and not re-proved here) or
\emph{physical inputs} (modeling assumptions outside mathematics).

\begin{abstract}
Classical spherical collapse reaches a central singularity~\cite{OppenheimerSnyder1939},
and singularity theorems convert trapped surfaces into incompleteness under energy
conditions~\cite{Penrose1965}.
Hysterical Response Theory (HRT) identifies the Hysterical force as the residual /
off-classical response of \emph{ordinary gravity} itself,
$\FH=-\Tr(F\mathcal L^{-1}\mathcal S)$, not a new carrier~\cite{boyd2026response}.
Gravity escapes Hysterical neutrality, and closed-loop response becomes unstable when
the gain crosses unity.
Here we apply that package to spherical collapse after horizon formation.
The single geometric input is causal failure of residual shell cancellation across the
horizon: the gravitational response left in the vacated exterior remains attractive
toward that wake, hence outward relative to the collapsing surface.
Identifying the wake acceleration with the radial HRT response, and invoking
HRT criticality under increasing curvature, yields an outward $H(R)$ that dominates
Newtonian self-gravity near $R=0$.
The reduced radial dynamics then admit a finite-radius steady remnant that is linearly
asymptotically stable under positive damping.
Thus HRT proves nonsingular spherical collapse in the reduced radial model: collapse
ends at $R_\ast>0$ rather than $R=0$.
\end{abstract}
\tableofcontents
\newpage

\section{Scope and status}
\label{sec:scope}
The response-theory paper defines Hysterical Response as the residual component of an
already-existing interaction under open-system dynamics~\cite{boyd2026response}:
\begin{equation}
\label{eq:FH}
\boxed{\FH:=-\Tr\!\bigl(F\,\mathcal L^{-1}\mathcal S\bigr).}
\end{equation}
No new force is added: $\FH$ is a component of the response of the underlying force
observable $F$.
For gravity, the Hysterical Neutrality Theorem does not force $\FH\to0$ (gravitational
survival), and the Neutrality--Instability Theorem supplies a supercritical phase when
the closed-loop gain crosses unity.

The present paper's claim is correspondingly strong:
\begin{equation}
\label{eq:claim}
\boxed{\begin{gathered}
\text{Hysterical Response Theory}
\;+\;
\text{spherical symmetry}
\;+\;
\text{horizon residual-cancellation failure}\\
\Longrightarrow\quad
\text{nonsingular spherical collapse (stable finite-radius remnant).}
\end{gathered}}
\end{equation}
Classical singularity theorems are not contradicted as pure GR statements; the dynamical
law for the gravitational residual after horizon formation is different once HRT is
included.

Four logical layers must be kept separate.
\begin{center}
\begin{tabular}{@{}>{\bfseries}p{0.28\linewidth}p{0.64\linewidth}@{}}
\toprule
Prior HRT theorem & $\FH$ definition; gravitational survival; neutrality--instability
(cited from~\cite{boyd2026response}).\\
Physical input & Horizon causes causal failure of residual shell cancellation.\\
Modeling assumption & Spherical reduction to areal radius $R(t)$; linear radial damping.\\
Theorem (this paper) & Wake identification; remnant existence/stability; nonsingular
collapse packaging.\\
\bottomrule
\end{tabular}
\end{center}

\section{Inherited Hysterical Response Theory}
\label{sec:hrt}
We work throughout in the gravitational channel of~\cite{boyd2026response}.

\begin{definition}[Gravitational Hysterical acceleration]
\label{def:gH}
Let $F$ be the gravitational force observable of the underlying Newtonian/GR-reduced
channel, $\mathcal L$ the stable open-system generator, and $\mathcal S$ an admissible
response source.
The Hysterical force is~\eqref{eq:FH}.
For a test mass $m_t$ at areal radius $R$, write
\begin{equation}
\label{eq:gH}
\mathbf g_H(R):=\frac{1}{m_t}\FH(R)
\end{equation}
for the corresponding acceleration (continuum packaging $\mathbf g_H=-\nabla\Phi_H$ as
in~\cite{boyd2026response}).
\end{definition}

\begin{proposition}[Gravitational survival; prior]
\label{prop:survival}
Ordinary positive-mass gravity is not response-visibly neutralized, so the Neutrality
Theorem does not force $\FH\to0$~\cite{boyd2026response}.
A nonzero gravitational Hysterical leftover is permitted.
\end{proposition}

\begin{proposition}[Hysterical instability; prior]
\label{prop:instab}
Under the closed-loop response hypotheses of~\cite{boyd2026response}, if the spectral
abscissa of the linearized gain Jacobian is positive (equivalently, reciprocal gain
exceeds unity), the neutral/controlled phase is linearly unstable.
In particular a closing Liouvillian gap $\gamma_0\to0^+$ with positive residue produces
a divergent Hysterical susceptibility and drives the channel through criticality.
\end{proposition}

\begin{remark}[Off-diagonal means residual]
\label{rem:offdiag}
``Off-diagonal'' / off-classical language in~\cite{boyd2026response} is pointer-relative:
$\FH$ is the residual response of ordinary gravity, not a separate interaction.
\end{remark}

\section{Horizon residual cancellation and the wake}
\label{sec:wake}
Classical spherical gravity obeys the shell theorem: exterior shells cancel for an
interior observer, and only interior mass sources $g(r)$~\cite{boyd2026response}.
The Hysterical residual is history-dependent and off-classical, so that cancellation is
an organizational property of the response, not an automatic identity once the matter
that maintained it has departed.

\begin{assumption}[Horizon residual-cancellation failure]
\label{ass:horizon}
Consider spherically symmetric matter of mass $M>0$ with areal radius $R(t)>0$.
After the matter crosses an event horizon, residual mutual cancellation among
gravitational shells fails causally for the Hysterical response channel:
the interior can no longer reorganize the exterior residual that was left in the region
previously occupied by the matter.
\end{assumption}

\begin{theorem}[Outward Hysterical wake]
\label{thm:wake}
Assume Definition~\ref{def:gH}, Proposition~\ref{prop:survival}, and
Assumption~\ref{ass:horizon}.
Let $\mathbf g_H^{\mathrm{wake}}$ be the gravitational Hysterical acceleration sourced by
the residual response remaining in the vacated exterior region.
Then, on the collapsing surface,
\begin{enumerate}
\item $\mathbf g_H^{\mathrm{wake}}$ is nonzero;
\item $\mathbf g_H^{\mathrm{wake}}$ is attractive toward the exterior wake;
\item relative to the outward normal of the collapsing sphere, the radial component
\begin{equation}
\label{eq:Hdef}
H(R):=\hat r\cdot\mathbf g_H^{\mathrm{wake}}(R)
\end{equation}
is a strictly positive outward acceleration.
\end{enumerate}
\end{theorem}

\begin{proof}
By Assumption~\ref{ass:horizon} the exterior residual is not cancelled by the
causally disconnected interior, so gravitational survival (Proposition~\ref{prop:survival})
applies to a nonzero leftover pairing: $\mathbf g_H^{\mathrm{wake}}\neq0$.
Attraction toward the wake is the sign of the underlying gravitational channel
(positive-mass residual response of ordinary gravity).
The wake lies outside the collapsing surface, so attraction toward it is radially
outward: $H(R)>0$.
\end{proof}

\begin{assumption}[Criticality under increasing curvature]
\label{ass:crit}
During continued collapse, increasing curvature drives the gravitational Hysterical
channel through the unstable phase of Proposition~\ref{prop:instab} (closing gap /
gain crossing unity).
Consequently the wake acceleration $H$ strengthens as $R$ decreases, and there exist
radii $0<R_{\mathrm{in}}<R_{\mathrm{out}}$ such that
\begin{equation}
\label{eq:sign}
H(R_{\mathrm{out}})<\frac{\GM}{R_{\mathrm{out}}^{2}}
\qquad\text{and}\qquad
H(R_{\mathrm{in}})>\frac{\GM}{R_{\mathrm{in}}^{2}},
\end{equation}
with $H$ continuous on $(0,\infty)$ and differentiable at equilibria used below.
\end{assumption}

\begin{remark}[What Ass.~\ref{ass:crit} uses from HRT]
\label{rem:crit}
The instability mechanism is Theorem material from~\cite{boyd2026response}.
That curvature supplies the concrete gap-closing for spherical collapse is the modeling
link stated in Assumption~\ref{ass:crit}.
\end{remark}

\section{Reduced radial dynamics}
\label{sec:radial}
\begin{model}[Spherical radial reduction]
\label{mod:radial}
Under spherical symmetry the surface dynamics reduce to the areal radius $R(t)>0$.
Net outward acceleration is
\begin{equation}
\label{eq:F}
F(R)\;:=\;H(R)-\frac{\GM}{R^{2}},
\end{equation}
with $H$ from Theorem~\ref{thm:wake} and $\GM/R^{2}$ the ordinary inward Newtonian
self-gravity.
The reduced radial dynamics are
\begin{equation}
\label{eq:ode}
\ddot R \;=\; F(R)\;-\;\gamma\,\dot R,
\end{equation}
with damping $\gamma\ge 0$ from dissipative open-system relaxation in the radial channel.
\end{model}

\begin{remark}[ODE status]
\label{rem:ode}
Equation~\eqref{eq:ode} is the spherical reduction of Newtonian self-gravity plus the
HRT wake acceleration of Theorem~\ref{thm:wake}, with linear damping.
It is not an independent force law: $H$ is gravitational $\FH$.
\end{remark}

\section{Existence of a remnant equilibrium}
\label{sec:existence}
\begin{lemma}[Intermediate-value zero]
\label{lem:ivt}
Let $a<b$ be positive reals and let $f:[a,b]\to\R$ be continuous.
If $f(a)>0$ and $f(b)<0$, then there exists $c\in(a,b)$ with $f(c)=0$.
\end{lemma}

\begin{proof}
Apply the intermediate value theorem to $f$ on $[a,b]$.
\end{proof}

\begin{theorem}[Existence of a finite-radius remnant]
\label{thm:existence}
Assume Theorem~\ref{thm:wake}, Assumption~\ref{ass:crit}, and Modeling
Assumption~\ref{mod:radial}.
Then there exists $R_\ast\in(R_{\mathrm{in}},R_{\mathrm{out}})$ such that
\begin{equation}
\label{eq:balance}
H(R_\ast)=\frac{\GM}{R_\ast^{2}},
\end{equation}
equivalently $F(R_\ast)=0$, and $R(t)\equiv R_\ast$ is a steady solution of~\eqref{eq:ode}.
\end{theorem}

\begin{proof}
By~\eqref{eq:sign}, $F(R_{\mathrm{in}})>0$ and $F(R_{\mathrm{out}})<0$ with $F$ continuous
on $[R_{\mathrm{in}},R_{\mathrm{out}}]$.
Lemma~\ref{lem:ivt} yields $R_\ast$ with $F(R_\ast)=0$.
\end{proof}

\section{Linear stability and dissipative attraction}
\label{sec:stability}
Linearising~\eqref{eq:ode} about $R_\ast$ gives
\begin{equation}
\label{eq:lin}
\delta\ddot R \;=\; F'(R_\ast)\,\delta R\;-\;\gamma\,\delta\dot R.
\end{equation}

\begin{lemma}[Restoring slope]
\label{lem:slope}
If $H'(R_\ast)<-2\GM/R_\ast^{3}$, then $F'(R_\ast)<0$.
\end{lemma}

\begin{proof}
$F'(R)=H'(R)+2\GM/R^{3}$.
\end{proof}

\begin{proposition}[Undamped restoring oscillation]
\label{prop:undamped}
If $\gamma=0$ and $F'(R_\ast)<0$, the linearisation is neutrally oscillatory with
frequency $\omega=\sqrt{-F'(R_\ast)}$.
\end{proposition}

\begin{theorem}[Damped attraction]
\label{thm:damped}
If $F'(R_\ast)<0$ and $\gamma>0$, every linear solution of~\eqref{eq:lin} decays to
zero as $t\to\infty$: all complex characteristic roots have negative real part, and the
standard underdamped, critically damped, and overdamped solution forms tend to zero.
\end{theorem}

\begin{proof}
Set $\kappa:=-F'(R_\ast)>0$.
The characteristic polynomial is $\lambda^{2}+\gamma\lambda+\kappa$.
Writing $z=x+iy$, the imaginary part forces $y(2x+\gamma)=0$; hence either a real root,
necessarily negative because sum $-\gamma<0$ and product $\kappa>0$, or
$x=-\gamma/2<0$.
Decay of the three standard real solution forms follows by elementary estimates
(envelope $e^{-(\gamma/2)t}$, or $e^{r t}$ with $r<0$).
\end{proof}

\begin{corollary}[No focusing to the origin in the linear basin]
\label{cor:nofocus}
Under Theorems~\ref{thm:existence} and~\ref{thm:damped}, linearised trajectories satisfy
$\delta R(t)\to0$, hence $R(t)\to R_\ast>0$.
\end{corollary}

\section{Nonsingular Gravitational Collapse Theorem}
\label{sec:main}
\begin{theorem}[Stable Hysterical Remnant]
\label{thm:remnant}
Under Theorem~\ref{thm:wake}, Assumption~\ref{ass:crit}, and Modeling
Assumption~\ref{mod:radial}:
\begin{enumerate}
\item there exists $R_\ast>0$ with $H(R_\ast)=\GM/R_\ast^{2}$;
\item if $H'(R_\ast)<-2\GM/R_\ast^{3}$ and $\gamma>0$, the remnant is asymptotically
stable for the linearised damped dynamics.
\end{enumerate}
\end{theorem}

\begin{proof}
Combine Theorems~\ref{thm:existence},~\ref{lem:slope}, and~\ref{thm:damped}.
\end{proof}

\begin{theorem}[Nonsingular Gravitational Collapse]
\label{thm:nonsingular}
Assume spherical symmetry, the gravitational HRT package of
Section~\ref{sec:hrt}~\cite{boyd2026response}, Assumption~\ref{ass:horizon},
Assumption~\ref{ass:crit}, Modeling Assumption~\ref{mod:radial}, the restoring slope
$H'(R_\ast)<-2\GM/R_\ast^{3}$, and $\gamma>0$.
Then spherical collapse does not terminate at a central shell-focusing singularity:
the reduced radial dynamics admit a stable finite-radius remnant $R_\ast>0$ to which
nearby linearised trajectories converge.
In that precise sense, Hysterical Response Theory proves nonsingular gravitational
collapse.
\end{theorem}

\begin{proof}
Theorem~\ref{thm:wake} supplies the outward HRT wake $H$.
Assumption~\ref{ass:crit} supplies the sign change~\eqref{eq:sign}.
Theorem~\ref{thm:remnant} then yields the stable remnant, and
Corollary~\ref{cor:nofocus} excludes linearised focusing to $R=0$.
\end{proof}

\begin{remark}[Scope of ``proves'']
\label{rem:proves}
The theorem is a necessity result inside HRT plus the stated geometric and reduction
hypotheses.
It does not claim a full Einstein--HRT PDE derivation of every coefficient, nor an
observational census of remnants.
The load-bearing physical input is Assumption~\ref{ass:horizon}; the instability engine
and the identification $H=\hat r\cdot\mathbf g_H$ are HRT.
\end{remark}

\section{Logical status of the results}
\label{sec:status}
\begin{center}
\small
\begin{tabular}{@{}p{0.48\linewidth}p{0.44\linewidth}@{}}
\toprule
\textbf{Quantity / claim} & \textbf{Status}\\
\midrule
$\FH=-\Tr(F\mathcal L^{-1}\mathcal S)$; off-classical residual & Prior HRT~\cite{boyd2026response}.\\
Gravitational survival; neutrality--instability & Prior HRT (Props.~\ref{prop:survival},~\ref{prop:instab}).\\
Horizon residual-cancellation failure & Physical input (Ass.~\ref{ass:horizon}).\\
Outward wake $H=\hat r\cdot\mathbf g_H^{\mathrm{wake}}$ & Theorem~\ref{thm:wake} (Lean).\\
Curvature $\Rightarrow$ critical wake dominance & Ass.~\ref{ass:crit} (HRT mechanism + collapse link).\\
Spherical radial ODE~\eqref{eq:ode} & Modeling Assumption~\ref{mod:radial}.\\
Remnant existence / damped stability & Theorems~\ref{thm:existence},~\ref{thm:damped} (Lean).\\
No-focus linear basin (all standard modes) & Corollary~\ref{cor:nofocus} (Lean).\\
Stable Hysterical Remnant & Theorem~\ref{thm:remnant} (Lean).\\
Nonsingular Gravitational Collapse & Theorem~\ref{thm:nonsingular} (Lean).\\
Full 4D Einstein--HRT PDE & Not claimed.\\
\bottomrule
\end{tabular}
\end{center}

\section{Conclusion}
\label{sec:conclusion}
Hysterical Response is the residual response of ordinary gravity.
After horizon formation, residual shell cancellation fails causally, leaving an exterior
gravitational wake whose HRT acceleration on the collapsing surface is outward.
HRT criticality under increasing curvature makes that wake dominate Newtonian
self-gravity near the classical singular limit.
The spherical radial reduction then necessarily admits a stable finite-radius remnant.

Thus, under the stated geometric and reduction hypotheses, Hysterical Response Theory
proves nonsingular spherical collapse: the endpoint is $R_\ast>0$, not $R=0$.

\bibliographystyle{plain}
\bibliography{paper_refs}
\end{document}
